\chapter*{ЗАКЛЮЧЕНИЕ}
\addcontentsline{toc}{chapter}{ЗАКЛЮЧЕНИЕ}

В ходе работы была достигнута поставленная во введении цель: разработано веб-приложение «Палантир», предназначенное для хранения, обработки и визуализации данных об исторических военных событиях.

В рамках решения первой задачи выполнен анализ предметной области, связанной с представлением пространственно-временных данных об исторических военных конфликтах. Установлено, что традиционные способы представления такой информации не позволяют в полной мере наглядно отразить динамику исторических военных событий. На основе проведённого анализа определены основные требования к разрабатываемому веб-приложению.

В соответствии со второй задачей рассмотрены технологические и архитектурные решения, необходимые для реализации веб-приложения. Сформирована архитектурная основа системы, предусматривающая клиентскую часть, серверное API, прикладные и предметные компоненты, слой доступа к данным и базу данных PostgreSQL/PostGIS.

В ходе решения третьей задачи разработана база данных и программная реализация веб-приложения «Палантир», включающие основные функции, связанные с хранением, обработкой и визуализацией данных. Рассмотрены структура базы данных, настройка Entity Framework Core, использование \texttt{PalantirDbContext}, организация репозиториев, взаимодействие слоёв приложения, подготовка серверного API, клиентской части и картографической визуализации пространственно-временных данных.

В рамках решения четвёртой задачи рассмотрены подходы к интеграционному тестированию, подготовке приложения к развёртыванию и проверке его работоспособности после размещения в серверной среде. Особое внимание уделено согласованности взаимодействия клиентской части, серверного API и базы данных, а также корректной передаче JSON- и GeoJSON-данных.

Практическая значимость выполненной работы заключается в создании веб-информационной системы, которая может служить основой для визуального представления данных об исторических военных событиях, а также для дальнейшего расширения функциональности и развития проекта.

Таким образом, цель работы достигнута, а задачи, сформулированные во введении, решены в полном объёме. Разработанная система объединяет базу данных, серверную логику, пользовательский интерфейс и средства картографической визуализации, что позволяет использовать её как основу для дальнейшего развития проекта «Палантир».